\documentclass[10pt,a4paper]{amsart}
\usepackage[margin=19mm]{geometry}
\usepackage{amsmath,amssymb,mathtools}
\usepackage[scr=rsfs,cal=euler]{mathalfa}
\usepackage{xfrac}
\usepackage[unicode,hidelinks,bookmarks=false]{hyperref}
\newcommand{\ie}{i.e.\ }
\newcommand{\cf}{cf.\ }
\newcommand{\resp}{respectively\ }
\newcommand{\Subsection}{\subsection}
\providecommand{\nobreakdash}{\nobreak\hbox{-}}
\setlength{\emergencystretch}{2em}
\multlinegap0pt
\newtheorem{lemma}{Lemma}
\newcommand{\PP}{\mathbb{P}}
\newcommand{\R}{\mathbb{R}}
\newcommand{\dK}{d_{\mathrm{Kol}}}
\newcommand{\eps}{\varepsilon}
\title{A Kolmogorov fourth-moment bound on Poisson chaos via a martingale core}
\author{Guangqu Zheng}
\date{Source: arXiv:2607.28742v1 (CC BY 4.0)}
\begin{document}
\maketitle

\noindent\emph{Source and licence.} A Kolmogorov fourth-moment bound on Poisson chaos via a martingale core. Version arXiv:2607.28742v1 (CC BY 4.0), distributed by arXiv under the Creative Commons Attribution 4.0 International licence, http://creativecommons.org/licenses/by/4.0/, as recorded in the arXiv licence metadata for this exact version and shown on the abstract page https://arxiv.org/abs/2607.28742v1. Full licence notice: \texttt{licenses/arxiv-2607.28742-CC-BY-4.0.txt}.\par\smallskip

\noindent\textit{Editorial scope.} Counted unit: the article's lemma lem\_lsc (source0.tex lines 1219--1227). The article's complete original proof of it is delivered with it (source0.tex lines 1229--1269, one \texttt{proof} environment of 41 lines). No mathematics is written, repaired or re-derived: the statement and the proof are the article's own lines, in the article's own notation, and no retained line is substituted or re-flowed.  No further source-local context is needed: every cross-reference of every retained line resolves inside the two delivered spans.  Nothing was omitted from any retained span, so the package carries no omission record.  No retained line contains a citation, so the wrapper carries no bibliography and no bibliography entry.  No retained line contains a figure environment, an image-inclusion command or drawing code, so no image is delivered and the package declares no asset dependency.  Editorial formatting only, in addition: an AMS \texttt{amsart} preamble that restates the article's own declarations and macros used by the retained lines, namely the environments \texttt{lemma} on separate counters, together with the standard packages those lines need.  The retained environments print in this document as Lemma 1 in order of appearance, whereas the article prints the same items with the numbering of its own full document.  The complete native source of the article as distributed in the arXiv e-print for this exact version is bundled unchanged as \texttt{sources/206436/source0.tex}.
\begin{lemma}[\textsf{Lower semicontinuity of Kolmogorov distance}]
\label{lem_lsc}
Let $X_n$, $X$, and $Y$ be real-valued random variables.  If
$X_n\to X$ in distribution and the distribution function of $Y$ is
continuous, then
$\dK(X,Y)
  \leq\liminf_{n\to\infty}\dK(X_n,Y).
$
\end{lemma}

\begin{proof}
Choose a subsequence $(n_k)$ along which
$
\dK(X_{n_k},Y)
\to 
\liminf_{n\to\infty} \dK(X_n,Y).
$
By   Skorohod representation theorem, this subsequence and $X$ may
be realized on a common probability space so that
$X_{n_k}\to X$ almost surely.  Let $F_Y$ be the distribution function
of $Y$ and define
$\omega_Y(\eps)
  :=\sup\{   (F_Y(x+\eps)-F_Y(x-\eps)) :x\in\R \}.
  $
A continuous distribution function is uniformly continuous on $\R$,
because it has finite limits at $-\infty$ and $+\infty$.  Hence
$\omega_Y(\eps)\to0$ as $\eps\downarrow0$.

For every $x\in\R$ and $\eps>0$,
$
  \PP(X\leq x)
  \leq\PP(X_{n_k}\leq x+\eps)
  +\PP(|X_{n_k}-X|>\eps)
$
and
$
  \PP(X_{n_k}\leq x-\eps)
  \leq\PP(X\leq x)
  +\PP(|X_{n_k}-X|>\eps).
$
Consequently,
\[
\begin{aligned}
  \bigl|\PP(X\leq x)-\PP(Y\leq x)\bigr|
  \leq{}&\dK(X_{n_k},Y)+\omega_Y(\eps)
  +\PP(|X_{n_k}-X|>\eps).
\end{aligned}
\]
Taking the supremum over $x$, letting $k\to\infty$, and then 
$\eps\downarrow0$, concludes the proof. 
\end{proof}

\end{document}
